% !TeX root = ../despliegue.tex
% ===========================================================================
%  Inventario de contenedores.
% ===========================================================================

\section{Qué se levanta}

Once contenedores, definidos en \cmd{docker-compose.yml}. Nueve se construyen
desde el propio repositorio —cada uno con su \cmd{Dockerfile} en
\cmd{frontend/<nombre>} o \cmd{backend/<nombre>}— y dos se descargan ya hechos:

\begin{center}
\small
\begin{xltabular}{\textwidth}{@{}l L{2.3cm} L{2.0cm} c Y@{}}
\toprule
\textbf{Servicio} & \textbf{Origen} & \textbf{Red} & \textbf{Puerto} &
\textbf{Para qué} \\
\midrule
\endfirsthead
\toprule
\textbf{Servicio} & \textbf{Origen} & \textbf{Red} & \textbf{Puerto} &
\textbf{Para qué} \\
\midrule
\endhead
\servicio{edge} & \cmd{nginx}\newline\cmd{1.27-alpine} & frontend\newline + backend & 80 &
  Único origen: sirve el shell, los remotes y \cmd{/api} \\
\servicio{shell} & build local & frontend & 3000 &
  Host de los microfrontends, sesión y enrutado \\
\servicio{mf-reservas} & build local & frontend & 3001 &
  Disponibilidad, nueva reserva, mis reservas \\
\servicio{mf-administracion} & build local & frontend & 3002 &
  Canchas, reservas y usuarios \\
\servicio{mf-reportes} & build local & frontend & 3003 &
  Los cuatro indicadores \\
\servicio{gateway} & build local & backend & 8080 &
  Verifica la credencial y rutea \\
\servicio{ms-usuarios} & build local & backend & 8081 &
  Registro, autenticación y usuarios \\
\servicio{ms-canchas} & build local & backend & 8082 &
  Catálogo, horarios y bloqueos \\
\servicio{ms-reservas} & build local & backend & 8083 &
  Reservas y disponibilidad \\
\servicio{ms-reportes} & build local & backend & 8084 &
  Indicadores; \textbf{sin base de datos} \\
\servicio{postgres} & \cmd{postgres}\newline\cmd{16-alpine} & backend & 5432 &
  Las tres bases, aisladas por credenciales \\
\bottomrule
\end{xltabular}
\end{center}

\textbf{El \servicio{edge} es el único contenedor conectado a las dos redes.} Es
lo que impide que el navegador alcance un microservicio saltándose el gateway:
los servicios de la red \emph{backend} sencillamente no son visibles desde
fuera.

Los puertos de los microservicios se publican solo para poder inspeccionar su
Swagger durante la evaluación; la aplicación no los usa.
